\documentclass[12pt]{article}
\usepackage[utf8]{inputenc}
\usepackage[T1]{fontenc}
\usepackage{hyperref}
\title{Teanu Language Resources for K-12 Curriculum Development}
\author{Compiled by Claude Sonnet 5 (claude-sonnet-5) for Andr\'as Kornai}
\date{}
\begin{document}
\maketitle

\section{Introduction}
Teanu (ISO 639-3 tkw; also called Puma or Buma) is an Oceanic language of Vanikoro, Temotu Province, Solomon Islands, with roughly 800 to 1000 speakers. It has replaced the neighbouring Lovono (now extinct) and Tanema (one speaker in 2012). The search found strong scholarly documentation by Alexandre Fran\c{c}ois, but no school materials, primers or curriculum resources.

\section{Primary Educational Resources}
No primers, textbooks or school programmes for Teanu turned up in the search results. The dictionary below could serve as a base for developing vocabulary materials.

\section{Linguistic and Reference Materials}
\begin{itemize}
\item Teanu--English dictionary (Dictionaria, A.\ Fran\c{c}ois, 2021): \url{https://dictionaria.clld.org/contributions/teanu}
\item Introduction to the Teanu dictionary (Fran\c{c}ois 2021, PDF): \url{https://marama.huma-num.fr/data/AlexFrancois_2021_Dictionaria_Teanu-dictionary_Introduction.pdf}
\item ``The languages of Vanikoro: Three lexicons and one grammar'' (Fran\c{c}ois 2009, HAL preprint): \url{https://halshs.archives-ouvertes.fr/halshs-00525557v1/file/AlexFrancois_2009_Vanikoro-languages_Festschrift-MDR_preprint.pdf}
\item The same paper as a PDF on the author's site: \url{https://marama.huma-num.fr/data/AlexFrancois_2009_Vanikoro-languages.pdf}
\item Author's Vanikoro notes page: \url{https://marama.huma-num.fr/AF-Vanikoro-notes_e.htm}
\item Verbix page on Teanu (verb list): \url{https://docs.verbix.com/Languages/Teanu}
\end{itemize}
A full Teanu reference grammar did not appear in the search results.

\section{Religious and Community Resources}
Nothing found in the search results.

\section{Digital and Audio Resources}
The Dictionaria entry is the only digital resource found. No audio collections were seen in the results.

\section{Curriculum Development Context}
Little exists for classroom use. Any curriculum would have to be built from the dictionary and linguistic papers above, ideally with Vanikoro community input. Context: \url{https://en.wikipedia.org/wiki/Languages_of_the_Solomon_Islands}
\end{document}